%------------------------------------------------------------------------------%
\begin{question}
  Diagonalize a matriz
  \[
    A=
    \begin{bmatrix}
      4 & 1 \\
      2 & 3
    \end{bmatrix}
  \]
  Determine \(P\) e \(D\) tais que
  \(
    A = PDP^{-1}
  \)
\end{question}

%------------------------------------------------------------------------------%
\begin{resolution}{Autovalores}
\begin{columns}
\column{0.4\textwidth}
  \onslide<+->{Polinômio característico}
  \begin{align*}
    \onslide<+->{p(\lambda)}
    \onslide<+->{& = \det(A-\lambda I)}
    \\
    \onslide<+->{& = \begin{vmatrix}[c]
      4-\lambda & 1         \\
      2         & 3-\lambda
    \end{vmatrix}}
    \\
  \onslide<+->{& = (4-\lambda)(3-\lambda) - 2 \\}
  \onslide<+->{& = 12 -7\lambda + \lambda^2 - 2 \\}
  \onslide<+->{& = \lambda^{2} - 7\lambda + 10}
  \end{align*}
\column{0.6\textwidth}
  \[
    \onslide<+->{\Delta}
    \onslide<+->{= b^2 - 4ac}
    \onslide<+->{= (-7)^2 - 4 \times 1 \times 10}
    \onslide<+->{= 9}
  \]

  \smallskip
  \[
    \onslide<+->{\sqrt{\Delta}}
    \onslide<+->{= \sqrt{9}}
    \onslide<+->{= 3}
  \]

  \smallskip
  \[
    \onslide<+->{\lambda_1}
    \onslide<+->{= \frac{-b+\sqrt{\Delta}}{2a}}
    \onslide<+->{= \frac{7+3}{2}}
    \onslide<+->{= 5}
  \]

  \smallskip
  \[
    \onslide<+->{\lambda_2}
    \onslide<+->{= \frac{-b-\sqrt{\Delta}}{2a}}
    \onslide<+->{= \frac{7-3}{2}}
    \onslide<+->{= 2}
  \]
\end{columns}
\end{resolution}

%------------------------------------------------------------------------------%
\begin{resolution}{Autovetores para $\lambda_1=5$}
\begin{columns}
\column{0.4\textwidth}
  \begin{align*}
    \onslide<+->{(A-5I)\vec{x}&=\vec{0}}
    \\
    \onslide<+->{\begin{bmatrix}[c]
      4 - 5& 1 \\
      2 & 3 - 5
    \end{bmatrix}
    \vetor{x}{y} & = \vetor{0}{0}}
    \\
    \onslide<+->{\begin{bmatrix}
      -1 &  1 \\
       2 & -2
    \end{bmatrix}
    \vetor{x}{y} & = \vetor{0}{0}}
  \end{align*}
  \begin{align*}
    \onslide<+->{-x + y & = 0} \\
    \onslide<+->{     y & = x}
  \end{align*}
\column{0.4\textwidth}
  \onslide<+->{Escolhendo $x=1$}
  \onslide<+->{\[
    y = 1
  \]}
  \onslide<+->{\[
    \vec{v}_1 = \vetor{1}{1}
  \]}
  \onslide<+->{\[
    E(5) = \ger \left\{ \vetor{1}{1} \right\}
  \]}
\end{columns}
\end{resolution}

%------------------------------------------------------------------------------%
\begin{resolution}{Autovetores para $\lambda_2=2$}
\begin{columns}
\column{0.4\textwidth}
  \begin{align*}
    \onslide<+->{(A-2I)\vec{x}&=\vec{0}}
    \\
    \onslide<+->{\begin{bmatrix}[c]
      4 - 2 & 1 \\
      2 & 3 - 2
    \end{bmatrix}
    \vetor{x}{y} & = \vetor{0}{0}}
    \\
    \onslide<+->{\begin{bmatrix}
      2 & 1 \\
      2 & 1
    \end{bmatrix}
    \vetor{x}{y} & = \vetor{0}{0}}
  \end{align*}
  \begin{align*}
    \onslide<+->{2x + y & = 0} \\
    \onslide<+->{     y & = -2x}
  \end{align*}
\column{0.4\textwidth}
  \onslide<+->{Escolhendo $x=1$}
  \onslide<+->{\[
    y = -2
  \]}
  \onslide<+->{\[
    \vec{v}_2 = \vetor{1}{-2}
  \]}
  \onslide<+->{\[
    E(2) = \ger \left\{ \vetor{1}{-2} \right\}
  \]}
\end{columns}
\end{resolution}

%------------------------------------------------------------------------------%
\begin{resolution}{Construção de \(P\) e \(D\)}
  Colocamos os autovetores nas colunas de \(P\)
  \pause
  \[
    P
    \pause
    = \left[ \vec{v}_2 \quad \vec{v}_2 \right]
    \pause
    = \begin{bmatrix}
      1 &  1  \\
      1 & -2
    \end{bmatrix}
  \]

  \pause
  Na mesma ordem dos autovalores de \(D\)
  \pause
  \[
    D
    \pause
    = \diag(\lambda_1, \lambda_2)
    \pause
    = \diag(5, 2)
    \pause
    = \begin{bmatrix}
        5 & 0 \\
        0 & 2
      \end{bmatrix}
  \]
\end{resolution}

%------------------------------------------------------------------------------%
\begin{proofframe}{Exemplo 1 -- Verificação -- Cálculo de \(P^{-1}\)}
  Para uma matriz de ordem \(2\)
  \[
    \begin{bmatrix}
      a & b \\
      c & d
    \end{bmatrix}^{-1}
    =
    \frac{1}{ad-bc}
    \begin{bmatrix}
      d  & -b \\
      -c & a
    \end{bmatrix}
  \]
  \pause
  Calculamos o determinante de \(P\)
  \pause
  \[
    \det(P)
    \pause
    = \begin{vmatrix}
      1 &  1  \\
      1 & -2
    \end{vmatrix}
    \pause
    = -2 -1
    \pause
    = -3
    \neq 0
  \]
  \pause
  \[
    P^{-1}
    \pause
    = \frac{1}{-3}
    \begin{bmatrix}
      -2 & -1 \\
      -1 & 1
    \end{bmatrix}
    \pause
    =
    \begin{bmatrix}
      \sfrac{2}{3} &  \sfrac{1}{3}  \\[3mm]
      \sfrac{1}{3} & -\sfrac{1}{3}
    \end{bmatrix}
  \]
\end{proofframe}

%------------------------------------------------------------------------------%
\begin{proofframe}{Exemplo 1 -- Verificação -- Diagonalização}
  \begin{align*}
    \onslide<+->{PDP^{-1}}
    \onslide<+->{& = \begin{bmatrix}
          1 & 1  \\
          1 & -2
        \end{bmatrix}
        \begin{bmatrix}
          5 & 0 \\
          0 & 2
        \end{bmatrix}
        P^{-1}}
    \\[2mm]
    \onslide<+->{& = \begin{bmatrix}
          5 & 2  \\
          5 & -4
        \end{bmatrix}
    P^{-1}}
    \\[2mm]
    \onslide<+->{& = \begin{bmatrix}
          5 & 2  \\
          5 & -4
        \end{bmatrix}
        \begin{bmatrix}
          \sfrac{2}{3} &  \sfrac{1}{3}  \\[3mm]
          \sfrac{1}{3} & -\sfrac{1}{3}
        \end{bmatrix}}
    \\[2mm]
    \onslide<+->{& = \begin{bmatrix}
          4 & 1 \\
          2 & 3
        \end{bmatrix}}
    \onslide<+->{= A}
  \end{align*}
\end{proofframe}

%------------------------------------------------------------------------------%
\endinput
